\section{Mean-ergodic theory and fixed-point structure}

\begin{definition}[Bounded-orbit endomorphism]
    \label{def:has_bounded_orbits}
    \lean{LinearMap.HasBoundedOrbits}
    \leanok
    Let $\mathbb K$ denote either $\R$ or $\C$, and let $V$ be a normed
    vector space over $\mathbb K$. A linear endomorphism $f:V\to V$ has
    \emph{bounded orbits} if, for every $x\in V$, the set
    $\{f^n x:n\in\N\}$ is bounded.
\end{definition}

\begin{lemma}[Ces\`aro averages of range elements]
    \label{lem:mean_ergodic_range_limit}
    \lean{LinearMap.HasBoundedOrbits.tendsto_birkhoffAverage_of_mem_range_sub_one}
    \leanok
    \uses{def:has_bounded_orbits}
    Let $f:V\to V$ have bounded orbits. If
    $x\in\range(f-\Id_V)$, then $A_N(x)\longrightarrow 0$.
\end{lemma}

\begin{proof}\leanok
    \uses{def:has_bounded_orbits}
    Write $x=(f-\Id_V)y$. Telescoping gives
    $A_N(x)=(f^Ny-y)/N$. Boundedness of the orbit of $y$ makes the
    numerator bounded, so the right-hand side tends to zero.
\end{proof}

\begin{lemma}[Complementarity of the fixed-point space]
    \label{lem:mean_ergodic_complementarity}
    \lean{LinearMap.HasBoundedOrbits.disjoint_ker_sub_one_range_sub_one,
        LinearMap.HasBoundedOrbits.isCompl_ker_sub_one_range_sub_one}
    \leanok
    \uses{def:has_bounded_orbits, lem:mean_ergodic_range_limit}
    Let $V$ be finite-dimensional and let $f:V\to V$ have bounded orbits.
    Then
    \begin{align}
        V
         & =\ker(f-\Id_V)\oplus\range(f-\Id_V).
        \label{eq:channel_mean_ergodic_decomp}
    \end{align}
\end{lemma}

\begin{proof}\leanok
    \uses{def:has_bounded_orbits, lem:mean_ergodic_range_limit}
    If $z$ belongs to both summands, then $A_N(z)=z$ because
    $z\in\ker(f-\Id_V)$, while Lemma~\ref{lem:mean_ergodic_range_limit} gives
    $A_N(z)\to0$. Hence $z=0$. Rank--nullity then gives
    (\ref{eq:channel_mean_ergodic_decomp}).
\end{proof}

\begin{definition}[Mean-ergodic projection]
    \label{def:mean_ergodic_projection}
    \lean{LinearMap.meanErgodicProjection}
    \leanok
    \uses{def:has_bounded_orbits, lem:mean_ergodic_complementarity}
    Let $V$ be finite-dimensional and let $f:V\to V$ have bounded orbits.
    The \emph{mean-ergodic projection} $P_f:V\to V$ is the projection onto
    $\ker(f-\Id_V)$ along $\range(f-\Id_V)$ in
    (\ref{eq:channel_mean_ergodic_decomp}).
\end{definition}

\begin{theorem}[Finite-dimensional mean-ergodic theorem]
    \label{thm:mean_ergodic_bounded_orbits}
    \lean{LinearMap.HasBoundedOrbits.tendsto_birkhoffAverage_meanErgodicProjection,
        LinearMap.HasBoundedOrbits.range_meanErgodicProjection,
        LinearMap.HasBoundedOrbits.meanErgodicProjection_apply_eq_self_iff,
        LinearMap.HasBoundedOrbits.isIdempotentElem_meanErgodicProjection,
        LinearMap.HasBoundedOrbits.meanErgodicProjection_apply_meanErgodicProjection,
        LinearMap.HasBoundedOrbits.comp_meanErgodicProjection,
        LinearMap.HasBoundedOrbits.meanErgodicProjection_comp}
    \leanok
    \uses{def:has_bounded_orbits, def:mean_ergodic_projection}
    Let $V$ be finite-dimensional and let $f:V\to V$ have bounded orbits.
    For every $x\in V$, the Ces\`aro averages satisfy
    \begin{align}
        A_N(x)
         & =\frac{1}{N}\sum_{n=0}^{N-1}f^nx
        \longrightarrow P_fx.
        \label{eq:channel_mean_ergodic_limit}
    \end{align}
    Moreover,
    \begin{align}
        \range(P_f) & =\ker(f-\Id_V), \notag \\
        P_f^2       & =P_f, \notag           \\
        fP_f        & =P_f=P_ff.
        \label{eq:channel_mean_ergodic_properties}
    \end{align}
    and $P_fx=x$ if and only if $fx=x$.
    The complex matrix specialization underlying
    \cite[Equation~(6.14)]{Wolf2012Quantum} is given below.
\end{theorem}

\begin{proof}\leanok
    \uses{def:has_bounded_orbits, lem:mean_ergodic_range_limit,
        lem:mean_ergodic_complementarity, def:mean_ergodic_projection}
    By Lemma~\ref{lem:mean_ergodic_complementarity}, decompose
    \begin{align}
        x      & =P_fx+(x-P_fx), \notag   \\
        P_fx   & \in\ker(f-\Id_V), \notag \\
        x-P_fx & \in\range(f-\Id_V).
    \end{align}
    The fixed component has constant Ces\`aro averages, whereas the averages
    of the second component tend to zero by
    Lemma~\ref{lem:mean_ergodic_range_limit}. This proves
    (\ref{eq:channel_mean_ergodic_limit}). Since
    $P_fx\in\ker(f-\Id_V)$, one has $f(P_fx)=P_fx$, and hence
    $fP_f=P_f$. Also, $(f-\Id_V)x\in\range(f-\Id_V)$ and
    $P_f((f-\Id_V)x)=0$, so $P_f(fx)=P_fx$ and hence $P_ff=P_f$.
    The range and idempotence identities in
    (\ref{eq:channel_mean_ergodic_properties}) follow directly from the
    projection construction.
\end{proof}

\begin{theorem}[Bounded orbits from trace nonincrease]
    \label{thm:positive_trace_nonincreasing_bounded_orbits}
    \lean{IsPositiveMap.hasBoundedOrbits_of_traceNonincreasing}
    \leanok
    \uses{def:positive_map, def:has_bounded_orbits}
    Let $T:\MN{D}\to\MN{D}$ be positive and trace nonincreasing. Then the
    forward orbit $\{T^n(X):n\geq0\}$ is bounded for every $X\in\MN{D}$.
    This is the trace-nonincreasing form of
    \cite[Proposition~6.3]{Wolf2012Quantum}.
\end{theorem}

\begin{proof}\leanok
    For $X\succeq0$, positivity and induction on $n$ give
    \begin{align}
        T^n(X)      & \succeq0, \notag \\
        \tr(T^n(X)) & \leq\tr(X).
    \end{align}
    Thus the orbit lies in a bounded trace section of the positive cone.
    Apply this to the positive and negative parts of a Hermitian matrix, and
    then write an arbitrary matrix as a complex linear combination of two
    Hermitian matrices.
\end{proof}

\begin{theorem}[Bounded orbits from unitality]
    \label{thm:positive_unital_bounded_orbits}
    \lean{IsPositiveMap.hasBoundedOrbits_of_unital}
    \leanok
    \uses{def:positive_map, def:unital_linear_map, def:has_bounded_orbits}
    Let $T:\MN{D}\to\MN{D}$ be positive and unital. Then the forward orbit
    $\{T^n(X):n\geq0\}$ is bounded for every $X\in\MN{D}$. This is the
    unital form of \cite[Proposition~6.2]{Wolf2012Quantum}, proof.
\end{theorem}

\begin{proof}\leanok
    For $X\succeq0$, positivity and unitality give, by induction on $n$,
    \[
        T^n(X)\preceq(\tr X)\cdot\Id,
    \]
    since $X\preceq(\tr X)\cdot\Id$ for $X\succeq0$ and $T$ is order-preserving
    with $T(\Id)=\Id$. Taking traces bounds $\tr(T^n(X))$ by $D\cdot\tr(X)$,
    so the orbit again lies in a bounded trace section of the positive cone.
    Apply this to the positive and negative parts of a Hermitian matrix, and
    then write an arbitrary matrix as a complex linear combination of two
    Hermitian matrices. This is the same uniform bound
    $\tr[A\,T(B)]\leq\|A\|_\infty\|B\|_\infty\tr[1\,T(1)]$ that gives
    Theorem~\ref{thm:positive_trace_nonincreasing_bounded_orbits} in the
    trace-nonincreasing case, specialized to $\tr[1\,T(1)]=\tr[1]=D$ via
    $T(1)=1$ instead of via trace preservation.
\end{proof}

\begin{theorem}[Positive trace-preserving mean-ergodic projection]
    \label{thm:positive_trace_preserving_mean_ergodic_projection}
    \lean{IsPositiveMap.hasBoundedOrbits_of_tracePreserving,
        IsPositiveMap.range_meanErgodicProjection_of_tracePreserving,
        IsPositiveMap.meanErgodicProjection_apply_eq_self_iff_of_tracePreserving}
    \leanok
    \uses{def:positive_map, def:trace_preserving,
        def:mean_ergodic_projection,
        thm:positive_mean_ergodic_projection,
        thm:trace_preserving_mean_ergodic_projection,
        thm:unital_mean_ergodic_projection}
    Let $T:\MN{D}\to\MN{D}$ be positive and trace-preserving. Then $T$ has
    bounded orbits, and its mean-ergodic projection $P_T$ is positive and
    trace-preserving. Its range is precisely the fixed-point space of $T$:
    \begin{align}
        \range(P_T) & =\ker(T-\Id), \notag \\
        P_T(X)=X    & \iff T(X)=X.
        \label{eq:channel_positive_mean_ergodic}
    \end{align}
    If $T(\Id)=\Id$, then $P_T(\Id)=\Id$.
    This is the Ces\`aro projection $T_\infty$ of
    \cite[Proposition~6.3 and Equation~(6.14)]{Wolf2012Quantum}.
\end{theorem}

\begin{proof}\leanok
    \uses{thm:mean_ergodic_bounded_orbits,
        thm:positive_trace_nonincreasing_bounded_orbits,
        thm:positive_mean_ergodic_projection,
        thm:trace_preserving_mean_ergodic_projection,
        thm:unital_mean_ergodic_projection}
    Trace preservation implies trace nonincrease, so
    Theorem~\ref{thm:positive_trace_nonincreasing_bounded_orbits} gives
    bounded orbits. Theorem~\ref{thm:mean_ergodic_bounded_orbits} then gives
    the projection and the fixed-point characterization
    (\ref{eq:channel_positive_mean_ergodic}). Its positivity follows from
    Theorem~\ref{thm:positive_mean_ergodic_projection}, trace preservation
    from Theorem~\ref{thm:trace_preserving_mean_ergodic_projection}, and its
    behavior on the identity from
    Theorem~\ref{thm:unital_mean_ergodic_projection}.
\end{proof}

\begin{theorem}[Positive unital adjoint Ces\`aro retraction]
    \label{thm:positive_unital_adjoint_mean_ergodic_retraction}
    \lean{IsPositiveMap.traceAdjoint_meanErgodicProjection_isPositiveUnitalRetraction}
    \leanok
    \uses{thm:positive_trace_preserving_mean_ergodic_projection,
        thm:trace_preserving_iff_trace_adjoint_unital,
        thm:adjoint_mean_ergodic_projection,
        thm:trace_pairing_adjoint_positive}
    Let $T:\MN{D}\to\MN{D}$ be positive and trace-preserving. The trace
    adjoint $P_T^*$ of its mean-ergodic projection is positive, unital, and
    idempotent. It is a retraction onto the adjoint fixed-point space:
    \begin{align}
        P_T^*(\Id)    & =\Id, \notag    \\
        (P_T^*)^2     & =P_T^*, \notag  \\
        \range(P_T^*) & =\ker(T^*-\Id).
        \label{eq:channel_adjoint_cesaro_retraction}
    \end{align}
    and $P_T^*(Y)=Y$ if and only if $T^*(Y)=Y$.
    This is the adjoint projection used in the proof of
    \cite[Theorem~6.14]{Wolf2012Quantum}.
\end{theorem}

\begin{proof}\leanok
    Theorem~\ref{thm:positive_trace_preserving_mean_ergodic_projection} makes
    $P_T$ positive and trace-preserving. Positivity passes to its trace
    adjoint, while
    Theorem~\ref{thm:trace_preserving_iff_trace_adjoint_unital} gives
    $P_T^*(\Id)=\Id$. The remaining assertions in
    (\ref{eq:channel_adjoint_cesaro_retraction}) are
    Theorem~\ref{thm:adjoint_mean_ergodic_projection}.
\end{proof}

\begin{theorem}[Positive fixed-point decomposition]\label{thm:positive_fixedpoint_parts}
    \lean{IsPositiveMap.posPart_negPart_fixed_of_fixedPoint}
    \leanok
    \uses{def:positive_map, def:trace_preserving}
    Let $E$ be a positive trace-preserving linear map and let $X=E(X)$ be a
    fixed point. Set $H_1=X+X^*$ and $H_2=i(X-X^*)$, and write
    $H_j=P_{j,+}-P_{j,-}$ for the canonical positive and negative parts.
    Then all four positive semidefinite matrices $P_{j,\pm}$ are fixed by
    $E$. This is \cite[Proposition~6.8]{Wolf2012Quantum}.
\end{theorem}

\begin{proof}
    \leanok
    Positivity implies that $E$ preserves adjoints. Hence $H_1$ and $H_2$
    are Hermitian fixed points. For either $H_j$, the fixed-point equation
    gives the common difference
    $Y=E(P_{j,+})-P_{j,+}=E(P_{j,-})-P_{j,-}$. Hence $P_{j,+}+Y$ and
    $P_{j,-}+Y$ are positive semidefinite, while trace preservation gives
    $\tr(Y)=0$. Diagonalize $P_{j,+}$. Orthogonality of the positive and
    negative parts shows that each eigenvector belongs to the kernel of one
    of them. Positivity of the corresponding matrix shows that every diagonal
    coefficient of $Y$ in this basis is non-negative. Their sum is zero, so
    every such coefficient vanishes. The positive semidefinite matrices then
    annihilate the corresponding basis vectors, and therefore $Y=0$.
    Applying this argument to $H_1$ and $H_2$ proves that all four parts are
    fixed.
\end{proof}
